% ============================================================
% Dịch vụ khử nhiễu (Noise Reduction Service)
% Dùng trong: Chương 3 (Thiết kế) và Chương 4 (Cài đặt)
% ============================================================

% --- THIẾT KẾ (được \input từ analysis-design.tex) ---

\newcommand{\designDenoise}{%

\section{Thiết kế dịch vụ khử nhiễu}

Dịch vụ khử nhiễu được thiết kế theo kiến trúc layered rõ ràng, tách biệt
các tầng xử lý để dễ bảo trì và mở rộng. Tầng API (\texttt{api/routes.py})
tiếp nhận HTTP request qua Flask Blueprint, thực hiện validate file đầu vào
và trả response. Tầng Service (\texttt{services/noise\_reduction.py}) chứa
logic nghiệp vụ: đọc file WAV, tiền xử lý audio (chuyển stereo sang mono,
normalize sang float32), gọi model inference, hậu xử lý kết quả (clip output,
ghi file). Tầng Model (\texttt{services/dtln\_model.py}) triển khai kiến trúc
DTLN bằng TensorFlow/Keras, xây dựng graph tính toán và load trọng số.

Dịch vụ cung cấp hai endpoint chính. Endpoint \texttt{POST /denoise} nhận
file WAV qua multipart/form-data (tối đa 50 MB), xử lý khử nhiễu và trả
về file WAV đã làm sạch. Endpoint \texttt{GET /health} trả về trạng thái
dịch vụ bao gồm: tình trạng model (ready/not ready), đường dẫn file trọng số,
sample rate, và input/output shape của model.

Mô hình DTLN được quản lý theo Singleton Pattern thông qua phương thức
\texttt{\_\_new\_\_()} --- model chỉ được load một lần duy nhất vào bộ nhớ
khi request đầu tiên đến, các request sau tái sử dụng cùng instance. Điều này
tránh chi phí khởi tạo lặp lại (mỗi lần load mất 2--3 giây) và đảm bảo
sử dụng bộ nhớ hiệu quả.

Framework được chọn là Flask 2.3+ kết hợp TensorFlow 2.13+. Flask phù hợp
với yêu cầu xử lý đồng bộ của pipeline khử nhiễu --- mỗi request cần chờ
model xử lý xong trước khi trả kết quả. CORS được bật cho tất cả origins
để cho phép gọi từ ESP32 và web client. Ứng dụng sử dụng Factory Pattern
(\texttt{create\_app()}) để khởi tạo, hỗ trợ cấu hình development và
production riêng biệt.

\begin{figure}[htbp]
\centering
\fbox{\parbox{0.85\linewidth}{\centering\vspace{2cm}
\textit{Thêm hình: Sơ đồ kiến trúc layered dịch vụ khử nhiễu --- 3 tầng:
API Layer (Flask Blueprint: /denoise, /health) $\rightarrow$ Service Layer
(NoiseReductionService: đọc WAV, tiền xử lý, hậu xử lý) $\rightarrow$
Model Layer (DTLN Singleton: STFT, LSTM, iFFT). Thể hiện luồng request
từ client qua từng tầng.}
\vspace{2cm}}}
\caption{Kiến trúc dịch vụ khử nhiễu}
\label{fig:denoise-architecture}
\end{figure}

}% end \designDenoise

% --- CÀI ĐẶT (được \input từ implementation.tex) ---

\newcommand{\implDenoise}{%

\section{Cài đặt dịch vụ khử nhiễu}

\subsection{Huấn luyện trên VIVOS}

Mô hình DTLN được huấn luyện trên bộ dữ liệu tiếng Việt VIVOS làm nguồn
giọng nói sạch, kết hợp với dữ liệu nhiễu Microsoft DNS ở các mức SNR
ngẫu nhiên để tạo cặp dữ liệu huấn luyện (audio nhiễu, audio sạch).
Quá trình huấn luyện sử dụng đoạn audio 3 giây (phù hợp với độ dài trung
bình 3--5 giây của VIVOS), batch size 32, và optimizer Adam với clipnorm 3.0
để tránh bùng nổ gradient. Learning rate ban đầu $10^{-3}$ được giảm tự
động qua ReduceLROnPlateau khi validation loss không cải thiện, kết hợp
Early Stopping sau 10 epoch không tiến triển để tránh overfitting.
\Cref{tab:training-config} tổng hợp cấu hình huấn luyện.

\begin{table}[htbp]
\centering
\caption{Cấu hình huấn luyện DTLN trên VIVOS}
\label{tab:training-config}
\begin{tabular}{>{\raggedright\arraybackslash}p{0.3\linewidth} >{\centering\arraybackslash}p{0.2\linewidth} >{\raggedright\arraybackslash}p{0.35\linewidth}}
\toprule
\rowcolor{headerblue} \textcolor{white}{\textbf{Tham số}} & \textcolor{white}{\textbf{Giá trị}} & \textcolor{white}{\textbf{Giải thích}} \\
\midrule
\rowcolor{rowgray} Dữ liệu sạch & VIVOS & Tiếng Việt, $>$ 1 GB \\
Dữ liệu nhiễu & Microsoft DNS & Kết hợp random SNR \\
\rowcolor{rowgray} Batch size & 32 & Cân bằng GPU/gradient \\
Độ dài đoạn & 3 giây & Phù hợp VIVOS (3--5s) \\
\rowcolor{rowgray} Epoch tối đa & 50 & Kết hợp dừng sớm \\
Learning rate & $10^{-3}$ & ReduceLROnPlateau \\
\rowcolor{rowgray} Early Stopping & 10 epoch & Tránh overfitting \\
Optimizer & Adam, clipnorm=3.0 & Tránh bùng nổ gradient \\
\bottomrule
\end{tabular}
\end{table}

Model sau huấn luyện được lưu dưới dạng file HDF5
(\texttt{DTLN\_vivos\_best.h5}, 3,9 MB) chứa toàn bộ trọng số của
khoảng 1,8 triệu tham số.

\subsection{Cài đặt mô hình DTLN}

Kiến trúc DTLN được triển khai trong file \texttt{dtln\_model.py} bằng
TensorFlow/Keras với các hyperparameter: 128 LSTM units, 2 lớp LSTM mỗi
giai đoạn, dropout 0.25, encoder size 256, block length 512 mẫu và
block shift 128 mẫu. Phương thức \texttt{build\_DTLN\_model()} xây dựng
graph tính toán gồm: lớp STFT (framing + rFFT), Separation Kernel thứ nhất
(2 LSTM + Dense + Sigmoid mask nhân với magnitude, kết hợp phase gốc, iFFT),
lớp Conv1D encoder (chiếu sang 256 chiều), Separation Kernel thứ hai
(2 LSTM + Dense + Sigmoid mask trong không gian đặc trưng), Conv1D decoder
(256$\rightarrow$512 với causal padding), và overlap-and-add để tái tạo
tín hiệu thời gian.

Lớp \texttt{InstantLayerNormalization} tùy chỉnh thực hiện chuẩn hóa
theo channel (axis=$-1$) với tham số gamma và beta có thể học được,
áp dụng công thức $(x - \mu) / \sqrt{\sigma^2 + \epsilon}$ với
$\epsilon = 10^{-7}$.

\subsection{Pipeline xử lý khử nhiễu}

Khi nhận request tại \texttt{POST /denoise}, tầng API thực hiện validate:
kiểm tra file tồn tại trong request, filename không rỗng, và extension là
\texttt{.wav} (duy nhất được chấp nhận). Filename được sanitize qua
\texttt{werkzeug.utils.secure\_filename} trước khi lưu tạm với prefix
\texttt{input\_} vào thư mục \texttt{uploads/}.

Tầng Service đọc file WAV qua thư viện SoundFile, tự động normalize sang
float32 trong phạm vi $[-1, 1]$. Nếu sample rate khác 16 kHz, hệ thống
ghi log cảnh báo nhưng vẫn tiếp tục xử lý. File stereo được convert sang
mono bằng cách lấy trung bình hai kênh (\texttt{np.mean(audio, axis=1)}).
Tín hiệu sau đó được thêm batch dimension (\texttt{np.expand\_dims}) và
đưa qua \texttt{model.predict(verbose=0)} để khử nhiễu. Đầu ra được
\texttt{np.clip} vào $[-1, 1]$ và ghi thành file WAV với prefix
\texttt{denoised\_}.

File kết quả được gửi về client qua \texttt{send\_file()} với MIME type
\texttt{audio/wav}. Cả file input và output tạm đều được xóa tự động
sau khi response gửi xong thông qua callback \texttt{response.call\_on\_close},
đảm bảo cleanup ngay cả khi có lỗi xảy ra trong quá trình gửi.

\subsection{Xử lý lỗi}

Dịch vụ xử lý 5 loại lỗi HTTP: 400 (file thiếu, filename rỗng, extension
không hợp lệ, audio format sai), 404 (file không tìm thấy), 413 (file
vượt quá 50 MB --- Flask tự xử lý qua \texttt{MAX\_CONTENT\_LENGTH}),
và 500 (lỗi nội bộ). Tất cả response lỗi trả về JSON với trường
\texttt{error} và \texttt{message}. File tạm luôn được cleanup kể cả
trong trường hợp exception.

}% end \implDenoise
